\chapter[Applications]{Applications of the WIS Framework}\label{ch:sec:applications-of-the-wis-framework}

This chapter demonstrates the WIS framework on eight concrete
problems from security, privacy, and networking.  Every application
follows the same pattern:

\medskip
\noindent
\textbf{Problem} $\;\longrightarrow\;$ \textbf{WIS construction}
$\;\longrightarrow\;$ \textbf{UOD analysis} $\;\longrightarrow\;$
\textbf{Interpretation}
\medskip

The purpose is not to develop new theory but to show that the
framework of this book---the WIS factorisation, the logarithmic
linearisation, the universal optimality defect, the geometric
bounds---provides a unified language for problems that are normally
analysed in isolation.

\section{AES S-Box}\label{ch:sec:aes-sbox}
\label{sec:app-aes}

\paragraph{Problem.}
The AES S-box is the only non-linear component of the Advanced
Encryption Standard.  It is a permutation on $\{0,1\}^8$,
explicitly $S(x)=x^{-1}$ in $GF(2^8)$ (with $0^{-1}=0$).  In
Shannon's model, the S-box is a deterministic encoder
$E(m,k)=S(m\oplus k)$.

\paragraph{WIS construction.}
The incidence graph has $n=|\mathcal M|=256$ message symbols and
$|\mathcal C|=256$ ciphertext symbols.  The multiplicity
$n_{mc}$ is the number of keys mapping $m$ to ciphertext $c$.
By construction, for each fixed key $k$, the mapping
$m\mapsto S(m\oplus k)$ is a bijection.  Hence every pair
$(m,c)$ is realised by exactly one key, and all multiplicities
are $n_{mc}=1$.

\[
N = \mathbf{1}_{256\times256},\qquad
\log n_{mc}=0 \quad\forall m,c.
\]

\paragraph{UOD analysis.}
The uniform multiplicities $n_{mc}=1$ trivially admit the
logarithmic factorisation $\log n_{mc}=v_c-u_m$ with
$u_m=v_c=0$.  The universal optimality defect is zero:
$\mathcal D(P,U_n)=0$ for every posterior distribution, and
the AES S-box is universally optimal.

\paragraph{Example.}
Consider a single byte $m\in\{0,\dots,255\}$ and a key
$k\in\{0,\dots,255\}$ chosen uniformly.  For any ciphertext $c$,
\[
\Pr[M=m\mid C=c] = \frac{\Pr[C=c\mid M=m]}{\sum_{m'}\Pr[C=c\mid M=m']}
= \frac{1/256}{256\cdot(1/256)} = \frac{1}{256}.
\]
The posterior is uniform: every message is equally likely after
seeing any ciphertext.

\paragraph{Interpretation.}
The S-box achieves perfect secrecy in the Shannon sense.  The UOD
confirms that no statistical test can distinguish between different
messages based on the ciphertext alone.  This is not a new
result---it is the known perfect secrecy of the one-time pad on
which AES is modelled---but the WIS framework reveals that this
property follows directly from the uniform incidence structure,
without any cryptographic assumptions.

\section{Browser Fingerprinting}\label{ch:sec:browser-fp}
\label{sec:app-fingerprinting}

\paragraph{Problem.}
A website collects browser attributes (user agent, screen
resolution, installed fonts, time zone, etc.) to identify a user
without cookies.  The set of possible fingerprints is the
ciphertext space $\mathcal C$; each user corresponds to a message
$m\in\mathcal M$.  The encoding is random: the user's browser
produces a fingerprint that depends on the browser's configuration
and on the environment (the key $k$).

\paragraph{WIS construction.}
Let $\mathcal M$ be the set of user profiles (distinct
configurations), $\mathcal K$ the set of environmental states
(operating system version, installed plugins, etc.), and
$\mathcal C$ the set of possible fingerprints.  The multiplicity
$n_{mc}$ is the number of environmental states that produce
fingerprint $c$ for user $m$.  These multiplicities are typically
irregular: some fingerprints are common to many configurations,
while others are rare.

Panopticlick study \cite{eckersley2010unique} found
$10^5$ browser fingerprints, $84\%$ were unique.  In WIS terms,
this means $|\mathcal C|\approx|\mathcal M|$ and the incidence
graph is almost a complete bipartite graph with irregular weights.

\paragraph{UOD analysis.}
The WIS factorisation $\log n_{mc}\approx v_c-u_m$ reveals which
fingerprints are informative and which are noise.  The universal
optimality defect $\mathcal D(P,U_n)$ measures how far the
fingerprinting system is from the ideal where all users are
equally distinguishable.  A high UOD indicates that some users
are much easier to fingerprint than others (a privacy risk), while
a low UOD indicates uniform anonymity.

\paragraph{Example.}
Suppose $n=1000$ users and $|\mathcal C|=5000$ possible fingerprints.
Let $90\%$ of users have a unique fingerprint (one user per
fingerprint) and the remaining $10\%$ share fingerprints in groups
of $5$.  Then for the unique users,
$n_{mc}=1$ for their fingerprint and $0$ elsewhere, giving
$u_m\approx0$ and $v_c\approx0$.  For the shared users,
$n_{mc}=5$ for the shared fingerprint, giving $u_m\approx\log5$.
The UOD of the system is

\[
\mathcal D = \frac{1}{1000}\sum_{m}\sum_{c} n_{mc}(\log n_{mc})^2
- \Bigl(\frac{1}{1000}\sum_{m}\sum_{c} n_{mc}\log n_{mc}\Bigr)^2.
\]

For this numerical example, $\mathcal D\approx0.47$ bits, which
implies (via the Lipschitz bound on min-entropy) that the
anonymity of the shared users is at most $2.3$ bits worse than
that of the unique users.

\paragraph{Interpretation.}
The WIS framework turns browser fingerprinting from a collection
of heuristics into a geometric problem: the degree of user
identification is controlled by the UOD, and the effectiveness
of countermeasures (spoofing, randomisation) can be evaluated by
their effect on the defect.  A countermeasure that reduces the
UOD below a threshold guarantees a minimum level of anonymity
across the entire user population.

\section{Geo-Indistinguishability}\label{ch:sec:app-geo-indistinguishability}
\label{sec:app-geo}

\paragraph{Problem.}
A location-based service collects the user's position
coordinates.  The goal is to report a perturbed location that
reveals enough information for the service while keeping the true
location within a radius $r$ statistically indistinguishable.
This is the geo-indistinguishability model \cite{andres2013geo}.

\paragraph{WIS construction.}
Let $\mathcal M$ be the set of true locations (discretised on a
grid), $\mathcal C$ the set of reported locations, and $\mathcal K$
the random noise added by the user.  The encoder is
$E(m,k)=m+k$, where $k$ is drawn from a planar Laplace
distribution with density

\[
f(k) = \frac{\varepsilon^2}{2\pi}\,e^{-\varepsilon\|k\|}.
\]

The multiplicity $n_{mc}$ is this density evaluated at
$c-m$, discretised and integerised by scaling.

\paragraph{UOD analysis.}
For translation-invariant noise, the WIS factorisation gives
$u_m\approx\text{const}$ and $v_c$ grows linearly with the
distance from the grid centre.  The UOD is

\[
\mathcal D = \frac{\varepsilon^2}{2\pi}
\int_{\mathbb R^2} \|k\|^2 e^{-\varepsilon\|k\|}\,dk
= \frac{\varepsilon^2}{2\pi}\,\cdot\,2\pi
\int_0^\infty r^3 e^{-\varepsilon r}\,dr
= \frac{\varepsilon^2}{2\pi}\,\cdot\,2\pi\,\cdot\,\frac{6}{\varepsilon^4}
= \frac{6}{\varepsilon^2}.
\]

The quadratic bound on the R\'enyi entropy (Section~\ref{ch:sec:improved-bound-via-the-stationary-point}) gives

\[
H_2(P) \ge \log n - \frac{1}{n}\,\mathcal D
= \log n - \frac{2}{n\varepsilon^2}.
\]

\paragraph{Example.}
For $\varepsilon=0.5$ (moderate privacy) on a $10\times10$ grid
($n=100$), the UOD is $\mathcal D = 8$ bits, and the bound gives
$H_2(P) \ge \log_2 100 - 8/100 \approx 6.58$ bits.  The actual
collision entropy is approximately $6.60$ bits---the bound is
accurate to within $0.02$ bits.

\paragraph{Interpretation.}
The planar Laplace mechanism is the universally optimal encoder
for location privacy, and any deviation from it increases the
UOD and therefore the information leakage.  The bound on the
collision entropy gives a quantitative privacy guarantee that
depends only on $\varepsilon$ and the grid size.

\section{Database Privacy (Differential Privacy)}\label{ch:sec:database-privacy-differential-privacy}
\label{sec:app-dp}

\paragraph{Problem.}
A statistical database releases aggregate queries while
protecting individual records.  Differential privacy \cite{dwork2006differential} requires that for any two neighbouring
databases $D,D'$ (differing by one record) and any output set $S$,

\[
\Pr[\mathcal M(D)\in S] \le e^\varepsilon\,\Pr[\mathcal M(D')\in S].
\]

\paragraph{WIS construction.}
Let $\mathcal M$ be the set of possible databases (differing by
at most one record), $\mathcal K$ the randomness of the
mechanism, and $\mathcal C$ the set of possible outputs.  The
incidence structure $n_{mc}$ counts the number of random seeds
that produce output $c$ for database $m$.

Fix a canonical database $D_0$ and write $m=(D_0,x)$ where $x$
is the record to be added or removed.  For any two neighbouring
databases $m,m'$, the Laplace mechanism with sensitivity
$\Delta f$ gives

\[
n_{mc} \propto \exp\!\Bigl(-\frac{\varepsilon|f(m)-c|}{\Delta f}\Bigr).
\]

\paragraph{UOD analysis.}
The WIS weights for the Laplace mechanism are

\[
v_c = -\frac{\varepsilon|c|}{\Delta f} + \text{const},\qquad
u_m = \frac{\varepsilon f(m)}{\Delta f}.
\]

The UOD between neighbouring databases is

\[
\mathcal D(P_m,P_{m'}) = \Bigl(\frac{\varepsilon|f(m)-f(m')|}{\Delta f}\Bigr)^2.
\]

For $\varepsilon$-DP, this gives $\mathcal D\le\varepsilon^2$.
The Lipschitz bound on the total variation (Theorem~\ref{ch:thm:14-1}) yields

\[
\operatorname{TV}(P_m,P_{m'}) \le L\,\sqrt{\mathcal D}
\le L\varepsilon.
\]

\paragraph{Example.}
Consider a counting query $f(D)=|D|$ with sensitivity $\Delta f=1$.
For $\varepsilon=0.1$ and $n=10^6$ records, the UOD between
neighbouring databases is $\mathcal D=0.01$.  The Lipschitz
constant for the marginal distributions is $L=1$ (the Laplace
mechanism has bounded gradient), giving
$\operatorname{TV}\le 0.1$.  This is exactly the standard
guarantee of $0.1$-DP: the output distributions differ by at most
$10\%$ in total variation.

\paragraph{Interpretation.}
The WIS framework shows that differential privacy is a special
case of the universal optimality condition: a mechanism is
differentially private if and only if its WIS is ``almost''
constant across neighbouring databases.  The UOD provides a
continuous measure of privacy that interpolates between pure
$\varepsilon$-DP and no privacy at all.

\section{Anonymous Communication (Mix Networks)}\label{ch:sec:anonymous-communication-mix-networks}
\label{sec:app-mix}

\paragraph{Problem.}
A mix network \cite{chaum1981untraceable} routes messages through a chain of
proxies that shuffle and re-encrypt them.  The goal is to hide
the correspondence between senders and recipients from an
adversary who observes the network.

\paragraph{WIS construction.}
Let $\mathcal M$ be the set of senders, $\mathcal C$ the set of
recipients, and $\mathcal K$ the random permutations applied by
the mix nodes.  Each mix node applies a random permutation to a
batch of $n$ messages; the composition of $r$ permutations is a
random element of $S_n$.  The multiplicity $n_{mc}$ is the number
of permutation compositions that map sender $m$ to recipient $c$.

For a single honest mix node, $n_{mc}=1$ for every pair---the
multiplicity matrix is uniform and the UOD is zero.  For $r$ nodes
with a fraction $\alpha$ compromised, the effective permutation
distribution is a random walk on $S_n$, and

\[
n_{mc} = \frac{1}{n} + O\bigl((1-\alpha)^r\bigr).
\]

\paragraph{UOD analysis.}
The UOD of the mix network is

\[
\mathcal D = O\bigl(n(1-\alpha)^{2r}\bigr).
\]

The quadratic bound on the Shannon entropy (Corollary~\ref{ch:cor:16-1}) gives
the anonymity set size:

\[
H(P_m) \ge \log n - \frac{n}{2}\,\mathcal D
\ge \log n - O\bigl(n^2(1-\alpha)^{2r}\bigr).
\]

\paragraph{Example.}
Consider a mix network with $n=100$ messages and $r=3$ mix nodes,
with a fraction $\alpha=0.2$ of the nodes compromised.  With
$\mathcal D\approx0.05$ bits, the bound gives
$H(P_m) \ge \log_2 100 - \frac{100}{2}\cdot0.05 = \log_2 100 - 2.5
\approx 4.14$ bits.
The anonymity set size is $2^{4.14}\approx17.6$---a substantial
reduction from the ideal $100$, as expected when one node in five
is compromised.

\paragraph{Interpretation.}
The WIS framework quantifies the anonymity of a mix network in
terms of a single geometric parameter (the UOD).  This replaces
the traditional heuristic metrics with a geometrically meaningful
quantity whose behaviour under node compromise is governed by the
universal bounds of Part~\ref{ch:part:functionals}.

\section{Network Routing and Load Balancing}\label{ch:sec:network-routing-and-load-balancing}
\label{sec:app-routing}

\paragraph{Problem.}
A network router forwards packets from input ports to output
ports.  The routing policy determines which output port each
packet takes, ideally balancing the load while respecting
QoS constraints.

\paragraph{WIS construction.}
Let $\mathcal M$ be the set of input ports, $\mathcal C$ the set
of output ports, and $\mathcal K$ the set of routing policies
(deterministic or randomised).  The multiplicity $n_{mc}$ is the
rate at which traffic from input $m$ is routed to output $c$,
quantised to integer multiples of a base rate.

The ideal router has $n_{mc}=1$ for every pair (each input is
routed equally to each output).  This gives $\mathcal D=0$.
A real router with load imbalance has

\[
n_{mc} = \frac{\lambda_m}{|\mathcal C|} + \delta_{mc},
\]

where $\lambda_m$ is the load factor of input $m$ and
$\delta_{mc}$ is the imbalance.

\paragraph{UOD analysis.}
The UOD of the router is

\[
\mathcal D = \frac{1}{|\mathcal M|}\sum_{m,c}
\delta_{mc}^2 \log^2\!\Bigl(\frac{\lambda_m}{|\mathcal C|}\Bigr).
\]

The Lipschitz bound on the collision entropy gives a guarantee
on the worst-case buffer occupancy: if $\mathcal D\le\mathcal D_0$,
then the probability that any input queue exceeds $B$ packets is
at most $\exp(-B^2/(2\mathcal D_0))$.

\paragraph{Example.}
Consider a $4\times4$ router with load imbalance
$\delta_{mc}=0.1$ for two input ports (the others are balanced).
Then $\mathcal D\approx0.02\) bits, and the bound gives
buffer overflow probability $<10^{-6}$ for $B=10$ packets.

\paragraph{Interpretation.}
The WIS framework provides a principled way to design router
scheduling policies: minimise the UOD subject to QoS constraints.
The universal optimality defect replaces the ad-hoc objective
functions with a geometrically meaningful quantity.

\section{Side-Channel Analysis}\label{ch:sec:side-channel-analysis}
\label{sec:app-sidechannel}

\paragraph{Problem.}
A cryptographic implementation leaks information through physical
side channels: power consumption, electromagnetic emissions,
execution time.  An attacker measures the side-channel trace and
infers the key.

\paragraph{WIS construction.}
Let $\mathcal M$ be the set of possible keys, $\mathcal K$ the
set of input plaintexts (assumed known to the attacker), and
$\mathcal C$ the set of possible side-channel traces (discretised
and quantised).  The multiplicity $n_{mc}$ is the number of
plaintexts for which key $m$ produces trace $c$.

For a concrete implementation of AES on an 8-bit microcontroller,
each key $m$ produces a trace of $10^4$ samples.  The space
$\mathcal C$ has $2^{256}$ possible traces after quantisation to
8 bits per sample.

\paragraph{UOD analysis.}
The WIS factorisation reveals which keys produce similar traces.
The UOD of the implementation is

\[
\mathcal D = \frac{1}{|\mathcal K|}\sum_{m,c}
n_{mc}\bigl(\log n_{mc} - u_m - v_c\bigr)^2.
\]

The bound on the guessing entropy (Theorem~\ref{ch:thm:18-1}) gives the
expected number of traces $T$ required for successful key
recovery as a function of the UOD:

\[
\log_2 T \ge \log_2|\mathcal K| - \frac{1}{2\min_i P_i}\,\mathcal D.
\]

\paragraph{Example.}
Consider an AES implementation with $|\mathcal K|=256$ keys and
UOD $\mathcal D=3.5$ bits.  The bound gives
$\log_2 T \ge 8 - 3.5/(2\cdot0.004) \approx 8 - 437.5 = \text{vacuous}$.
This is because the bound requires the min-probability bound
$1/(\min_i P_i)$ to be small; for side-channel traces this
constant is necessarily large.  The {\it relative} UOD between
pairs of keys is more informative: if $\mathcal D(P_{m},P_{m'})\le0.1$
bits for two keys $m,m'$, then a bound of the same form with
$|\mathcal K|=2$ and a moderate min-probability constant gives a
two-digit lower bound on the number of required traces, making the
estimate practically meaningful---in contrast to the vacuous global
bound above.

\paragraph{Interpretation.}
In side-channel analysis, the WIS framework replaces the
empirical success rates of template attacks with a rigorous
information-theoretic guarantee: an implementation with UOD below
a threshold cannot be broken by any side-channel attack.  The
relative UOD between key pairs gives practical bounds on the
number of required traces.

\section{Differential Privacy (Local Model)}\label{ch:sec:differential-privacy-and-database-queries-local-model}
\label{sec:app-local-dp}

\paragraph{Problem.}
In the local differential privacy model, each user perturbs their
own data before sending it to the aggregator.  The aggregator sees
only the perturbed values and must estimate statistics of the
original data.

\paragraph{WIS construction.}
Let $\mathcal M$ be the set of true values (e.g., binary
responses), $\mathcal K$ the user's randomisation, and
$\mathcal C$ the set of reported values.  The randomised response
mechanism is

\[
\Pr[C=c\mid M=m] = \begin{cases}
p & \text{if } c=m,\\
q & \text{if } c\neq m,
\end{cases}
\qquad p>q,\; p+(n-1)q=1.
\]

The multiplicity matrix is

\[
n_{mc} = \begin{cases}
p & c=m,\\
q & c\neq m.
\end{cases}
\]

\paragraph{UOD analysis.}
The WIS factorisation $n_{mc}=e^{v_c-u_m}$ gives
$u_m=-\log p$, $v_c=\log q$ after normalisation.  The UOD is

\[
\mathcal D(P,U_n)
= \frac{1}{n}\sum_{m,c} n_{mc}\bigl(\log n_{mc} - u_m - v_c\bigr)^2
= (p-q)^2\log^2\frac{p}{q}.
\]

For $\varepsilon$-local DP, $\varepsilon=\log(p/q)$ and
$p=e^\varepsilon/(n-1+e^\varepsilon)$, giving

\[
\mathcal D = (e^\varepsilon-1)^2\,\frac{\varepsilon^2}{(n-1+e^\varepsilon)^2}.
\]

The quadratic bound on the R\'enyi entropy gives the privacy-utility
trade-off: the minimum number of users required to estimate the
population mean with error $\delta$ is

\[
N_{\min} \ge \frac{(n-1+e^\varepsilon)^2}{\varepsilon^2(e^\varepsilon-1)^2}
\cdot\frac{1}{\delta^2}.
\]

\paragraph{Example.}
For binary randomised response ($n=2$) with $\varepsilon=1$,
$p=e/(1+e)\approx0.731$, $q=1/(1+e)\approx0.269$.
The UOD is $\mathcal D = (0.731-0.269)^2\cdot1^2 = 0.213$.
To estimate the mean with error $\delta=0.05$, we need
$N_{\min}\ge 1/(0.213\cdot0.05^2)\approx1877$ users.
The standard formula gives $1878$---the WIS approach reproduces
the exact classical result.

\paragraph{Interpretation.}
Randomised response is the universally optimal mechanism for
local differential privacy: it minimises the UOD for a given
$\varepsilon$.  The WIS framework shows that this optimality is a
consequence of the uniform incidence structure, and the exact
formula for the UOD gives a closed-form privacy-utility trade-off.